% 星座（综述）
% license CCBYSA3
% type Wiki

本文根据 CC-BY-SA 协议转载翻译自维基百科\href{https://en.wikipedia.org/wiki/Constellation}{相关文章}。

\begin{figure}[ht]
\centering
\includegraphics[width=6cm]{./figures/a943f80cad7e3c1f.png}
\caption{} \label{fig_Coti_1}
\end{figure}
星座是天球上的一个区域，在该区域内，一组可见的恒星形成了人们所感知的图案或轮廓，通常代表一种动物、神话题材，或无生命的物体。$^{[1]}$

最早的星座很可能是在史前时期被定义的。人们利用星座来讲述关于其信仰、经历、创世观念以及神话的故事。不同的文化和国家创造了各自的星座体系，其中一些一直延续到20世纪初，直到当今的星座体系在国际上获得统一认可。随着时间的推移，对星座的认知发生了显著变化。许多星座的大小或形状发生了改变。有些一度流行，后来却逐渐被遗忘。有些则仅局限于单一文化或国家。对星座进行命名也帮助天文学家和航海者更容易地识别恒星。$^{[2]}$

十二个（或十三个）古代星座属于黄道星座，它们横跨黄道——太阳、月球和行星运行所经过的路径。黄道的起源在历史上仍不确定；其占星学上的划分大约在公元前400年于巴比伦或迦勒底天文学中变得突出。$^{[3]}$ 星座通过希腊文化进入西方世界，并出现在赫西俄德、欧多克索斯和阿拉托斯的著作中。由黄道星座及另外36个星座组成的传统48星座（由于阿尔戈船座被划分为三个星座，如今为38个），由出生于埃及亚历山大的希腊—罗马天文学家托勒密在其著作《天文学大成》中列出。星座的形成是大量神话的主题，其中最著名的当属拉丁诗人奥维德的《变形记》。15世纪至18世纪中叶，随着欧洲探险者开始前往南半球，远南天区的星座被陆续加入。由于罗马和欧洲文化的传播，每个星座都拥有一个拉丁名称。

1922年，国际天文学联合会正式接受了现代的88个星座列表，并于1928年采用了官方的星座边界，这些边界共同覆盖了整个天球。$^{[4][5]}$ 在天球坐标系统中，任意一个给定点都位于某一个现代星座之内。一些天文学命名体系在命名天体时会包含其所在的星座，以传达该天体在天空中的大致位置。例如，恒星的弗兰斯蒂德命名法由一个数字和星座名称的所有格形式组成。

其他被称为星群的恒星图案或组合，在严格定义下并不属于星座，但同样被观测者用来辨认夜空。星群可能是某一星座中的几颗恒星，也可能与多个星座共享恒星。星群的例子包括人马座中的茶壶，以及大熊座中的北斗七星。$^{[6][7]}$
\subsection{术语}  
“星座”一词源自晚期拉丁语 cōnstellātiō，可译为“恒星的集合”；该词于14世纪进入中古英语。$^{[8]}$ 古希腊语中表示星座的词是 ἄστρον（astron）。这些术语在历史上泛指任何可被识别的恒星图案，其外观通常与神话人物或生物、地面动物或某种物体相关联。$^{[1]}$ 随着时间推移，在欧洲天文学家之间，星座逐渐被明确界定并获得广泛认可。20世纪，国际天文学联合会（IAU）正式确认了88个星座。$^{[9]}$

当从地球上某一特定纬度观测时，若某个星座或恒星永远不会落到地平线以下，则称其为拱极星座。从北极或南极观测，位于天赤道以南或以北的所有星座都是拱极的。根据不同定义，赤道星座可以指位于北纬45°至南纬45°之间的星座$^{[10]}$，也可以指穿过黄道（或黄道带）赤纬范围、即北纬23.5°至南纬23.5°之间的星座。$^{[11][12]}$

星座中的恒星在天空中看似彼此接近，但它们通常位于距离地球各不相同的位置。由于每颗恒星都具有各自独立的运动，所有星座都会随着时间缓慢发生变化。经过数万至数十万年，熟悉的星座轮廓将变得难以辨认。$^{[13]}$天文学家可以通过精确测量单个恒星的共同自行$^{[14]}$（依赖高精度天体测量学$^{[15][16]}$）以及通过天文光谱学测定其径向速度$^{[17]}$，来预测星座在过去或未来的形态。

国际天文学联合会所承认的88个星座，以及历史上各文化所使用的星座，都是基于可观测天空中恒星分布而想象出的形象与形状。许多官方承认的星座源自古代近东和地中海地区的神话想象。$^{[18][19][page needed]}$ 其中一些故事似乎与星座的外观有关，例如猎户座被天蝎座刺杀的传说，其对应的星座在一年中的不同季节出现在天空中。$^{[20]}$
\subsection{观测}
\begin{figure}[ht]
\centering
\includegraphics[width=8cm]{./figures/479bf5216e7afbe4.png}
\caption{} \label{fig_Coti_2}
\end{figure}
由于地球围绕太阳运行的轨道上，夜晚发生的时间对应于轨道上的位置不断变化，星座在一年之中其位置也会随之改变。随着地球自转方向指向东方，天球看起来仿佛向西旋转，恒星在北极星周围呈逆时针环绕，在南极星周围则呈顺时针环绕。$^{[21]}$

由于地球自转轴具有23.5°的倾角，黄道沿着黄道面在南、北半球中均匀分布，近似构成一个大圆。北天的黄道星座包括双鱼座、白羊座、金牛座、双子座、巨蟹座和狮子座；南天的黄道星座包括室女座、天秤座、天蝎座、人马座、摩羯座和宝瓶座。$^{[22][a]}$ 根据一年中的不同时间，从北纬23.5°到南纬23.5°的纬度范围内，黄道会出现在头顶正上方。夏季时，黄道在白天显得较高、夜晚较低；而在冬季，情况则正好相反，这一现象在南、北半球均一致。

由于太阳系相对于银河系具有约60°的倾角，银河系的银道面相对于黄道倾斜约60°，位于北天的金牛座与双子座之间，以及南天的天蝎座与人马座之间，银河系中心也位于这一南天区域附近。$^{[22][23]}$ 银河在天空中穿过天鹰座（靠近天赤道），以及北天的天鹅座、仙后座、英仙座、御夫座和猎户座（靠近参宿四），同时也穿过麒麟座（靠近天赤道），以及南天的船尾座、船帆座、船底座、南十字座、半人马座、南三角座和天坛座。$^{[22]}$
\subsubsection{北半球}  
北极星作为北极星，是北天球的近似中心。它属于小熊座，位于小北斗勺柄的末端。$^{[22]}$

在大约北纬35°的地区，一月时，大熊座（包含北斗七星）出现在东北方，而仙后座位于西北方。西方可见双鱼座（位于地平线上方）和白羊座。西南方的鲸鱼座接近地平线。高空偏南的位置可见猎户座和金牛座。东南方向地平线上方是大犬座。在猎户座的上方偏东方向可以看到双子座；东方方向（并逐渐接近地平线）还有巨蟹座和狮子座。除金牛座外，英仙座和御夫座也出现在头顶附近。$^{[22]}$

在相同纬度下，七月时，仙后座（位于低空）和仙王座出现在东北方。大熊座此时位于西北方。牧夫座高悬于西方天空。室女座位于西方，天秤座在西南方，天蝎座在南方。人马座和摩羯座位于东南方。天鹅座（包含北十字星）位于东方。武仙座与北冕座一同高悬于天空中。$^{[22]}$


